\section{Shared Experimental Details}
\label{app:setup}

Sections~\ref{sec:normalization}--\ref{sec:density} contain the experiments;
this appendix records shared resource and evaluation details.

\paragraph{Training setup and run reuse.}
Section~\ref{sec:normalization} specifies the student, domains, and three
teacher weight sets. The $64$-response update comprises $16$ micro-batches
of four responses, with equal domain quotas, one optimizer update per
fresh rollout batch, and a $4{,}096$-token cap.
The initial budget is at most $500$ updates per run, subject to a
throughput pilot on a $96$-GB learner GPU and $48$-GB rollout/teacher GPUs.
Exact model revisions, optimizer and decoding settings, data splits,
and evaluation versions remain to be frozen before execution.

A representative single-teacher setting is chosen before inspecting
results. PG and corrected teacher top-$64$ in the single-teacher and
joint settings give four configurations. The joint top-$64$ configuration
with domain-response averaging is also the reference for the three-way
normalization study, which adds two averaging rules. Thus six training
configurations cover the core comparisons, with checkpoints reused for
the gradient, teacher-overlap, and full-vocabulary measurements.
Additional teacher settings, the response-cap comparison, and repeated
seeds require separate resource accounting. The manuscript does not yet
report measured outcomes for these studies.

\paragraph{Capability evaluation.}
Use MATH-500 greedy pass@1, a frozen LiveCodeBench pass@1 slice,
IFBench strict accuracy, and GPQA-Diamond average@4. Record evaluator,
model, and data revisions and fix decoding within each comparison.
Report raw per-domain scores, their macro-average, and the worst-domain
change from initialization. Independently trained OPD students, where
available, are references at matched domain exposure, not upper bounds
or guarantees of jointly attainable capability. A lack of single-teacher
transfer is not itself evidence of multi-teacher interference.
Loss on a common held-out response bank and loss on fresh student
responses are reported separately; a lower mean KL alone does not
establish balanced capability integration.

\paragraph{Compute and uncertainty.}
Record generated tokens, domain response counts, optimizer updates, and
GPU hours separately, including teacher-scoring and diagnostic overhead.
Key online comparisons use repeated paired training seeds when feasible;
initial one-seed runs are exploratory. Prompt bootstrap estimates
conditional evaluation uncertainty, not training-seed variation.
Report effect sizes, confidence intervals, and the range of differences
the experiment can resolve. Teacher pairs sharing weights or training
lineage, and repeated checkpoints, are not independent observations.
No sparsity advantage or capability improvement is claimed before
measurement.
